\documentclass[11pt]{article}
\usepackage[a4paper,margin=24mm]{geometry}
\usepackage[T1]{fontenc}
\usepackage[utf8]{inputenc}
\usepackage{lmodern,microtype}
\usepackage{amsmath,amssymb,amsthm}
\usepackage{booktabs,longtable,array,tabularx}
\usepackage{enumitem}
\usepackage{xcolor}
\usepackage[hidelinks]{hyperref}
\hypersetup{pdftitle={Six attempts at quasi-polynomial unknot recognition: synthesis, cross-validation, and a faster exact recognizer}}
\setlength{\parskip}{4pt}
\setlength{\parindent}{0pt}
\setlist{nosep,leftmargin=*}
\newtheorem{proposition}{Proposition}[section]
\newtheorem{lemma}[proposition]{Lemma}
\theoremstyle{remark}
\newtheorem{remark}[proposition]{Remark}
\newcommand{\F}{\mathbb{F}_2}
\newcommand{\Kh}{\widetilde{\mathrm{Kh}}}
\newcommand{\code}[1]{\texttt{#1}}
\newcommand{\qp}{n^{O(\log n)}}

\title{Six attempts at quasi-polynomial unknot recognition:\\
synthesis, cross-validation, and a faster exact recognizer}
\author{Synthesis report prepared with Claude Code}
\date{17 September 2026}

\begin{document}
\maketitle

\begin{abstract}
Six independently produced archives each tried to implement the $\qp$ unknot
recognition algorithm announced in Marc Lackenby's February 2021 talk. This
report extracts them into numbered directories, reads their reports, reviews
their code, re-runs their test suites, and cross-validates them against each
other on shared inputs. All six reach the same honest conclusion: the
quasi-polynomial algorithm was not implemented, because the accelerated
hierarchy operations named on the last slide have no published algorithmic
specification with cost bounds. What they deliver instead is consistent and
correct as far as we can test: five reduced-Khovanov recognizers agree on 44
shared inputs, six ball-pattern testers agree on 2000 random spherical
patterns, and the grid-diagram search of archive 01 agrees with Khovanov
homology on every rectangular diagram we checked. Building on their ideas we
add a seventh recognizer, \code{fastunknot}, whose exact part is a scanning
computation of Khovanov homology (Bar-Natan's local algorithm with delooping
and Gaussian elimination) preceded by an Alexander-polynomial filter. It is
still exponential in the worst case, but its cost is governed by the width of
a scan through the diagram rather than by the crossing number: the
80-crossing unknot family on which every cube implementation needs $2^{80}$
states is decided in eleven milliseconds. The $\qp$ target remains open; the
final sections state precisely what is missing.
\end{abstract}

\tableofcontents

\section{Task, sources, and outcome}

\subsection{What was asked}
The original request to each of the six runs was an implementation of the
main theorem of the supplied talk \cite{slides}: an algorithm that decides
whether an $n$-crossing knot diagram represents the unknot in time $\qp$.
The present task was to unpack the six resulting archives, read their reports,
review their code, and synthesize what was achieved and what obstacles remain.
A second task was added mid-way: write the asymptotically fastest algorithm
possible, trying to reach $\qp$.

\subsection{The talk}
The supplied file \code{docs/quasipolynomial-talk.pdf} has 109 pages including
progressive-reveal duplicates. Its logical content is short:
\begin{itemize}
\item pp.~11--14: the theorem, $\qp$ for an $n$-crossing diagram;
\item pp.~15--30: hierarchies, boundary patterns, essential patterns (a proper
  disc meeting the pattern at most three times must cut off an empty disc, an
  arc, or a tripod), and the Waldhausen--Johannson characterization: the knot
  is nontrivial iff the exterior admits an essential hierarchy ending in balls;
\item pp.~38--53: simplification of an inessential hierarchy by a violating disc
  (compress or pattern-compress the last surface the disc runs over, discard
  later surfaces), and the basic algorithm as a flowchart;
\item pp.~54--65: termination by a lexicographic complexity $(g(S_1),\dots,g(S_\ell))$,
  with the warning that ``genus'' may mean $\chi_-$ or a pattern-sensitive
  variant, and the estimate ``running time at most $Lg^L$'' if every surface
  has complexity at most $g$ and hierarchies have length at most $L$;
\item pp.~66--74: the four speed-ups needed for $\qp$: (1) surfaces with
  $g\le O(n^2)$, (2) compressed hierarchy encodings (Agol--Hass--Thurston),
  (3) multi-surfaces, (4) generalised Heegaard splittings and Cheeger regions,
  ``the effect of this is to ensure that $L\le O(\log n)$'';
\item pp.~75--108: sketches of each speed-up, including the Cheeger-region
  condition $g(\partial M_i)\le\frac13\min\{g(H_i),g(H)-g(H_i)\}$;
\item p.~109: the final flowchart with the boxes \textsc{Build hierarchical
  multi-surface}, \textsc{Cut along surface}, \textsc{Use Cheeger region},
  \textsc{Haken's lemma}, \textsc{Weakly reduce}, \textsc{Simplify
  multi-surface}, and the terminal test ``does the graph dual to the boundary
  pattern have a violating cycle?''.
\end{itemize}
The slides describe the strategy; they contain no data structures, no
algorithms for the boxes, and no proof of the depth bound.

\subsection{The 2026 preprint}
All six archives located Lackenby's July 2026 preprint \cite{lackenby2026}
(``Incompressible surfaces, hierarchies and unknot recognition'', 54 pages).
We downloaded it and read Section 9 (``The number of steps'') ourselves. It
opens: ``There are various steps in the algorithm from Theorem 1.3 that are
not described with enough precision to be able to estimate their running
time. So, we will initially content ourselves with an analysis of the number
of times that the algorithm goes through any of the steps'', in terms of the
maximal hierarchy length $L$ and the maximal pattern-complexity $g$, and ``At
present, we will not attempt to bound these two quantities.'' Proposition 9.1
then gives the iteration bound $L(g+1)^L$ using exactly the base-$(g+1)$
digit argument that the talk sketches in base $g$. Section 10 refines the
algorithm to obtain a bound on $L$ that is \emph{linear} in the initial
$q$-complexity of the handle structure, not logarithmic in $n$. The preprint
also supplies two facts the archives use: Proposition 2.6 (a pattern on a
3-ball is essential iff it is empty, a circle, or a connected graph whose
planar dual is simple and 4-connected or $K_3$ or $K_4$) and Definition 3.4
(pattern-complexity $\chi_P(S)=-4\chi(S)+|S\cap P|$).

So the accelerated, logarithmic-depth algorithm of the talk is, as of
September 2026, announced but not specified in any source available to us.
This is the single obstacle behind every ``partial'' verdict below, and it is
mathematical, not an engineering shortfall.

\subsection{Outcome in one paragraph}
None of the six archives, and not the new package either, implements a
recognizer with a proved $\qp$ bound. Each archive contains a complete exact
recognizer with an exponential bound, plus genuinely implemented polynomial
pieces of the hierarchy strategy (the terminal ball-pattern test, cocycle-dual
normal coordinates, the potential arithmetic), and each report says so
without hedging. We verified every test suite, cross-validated the archives
against each other and against independent computations, found no incorrect
verdict, and identified one gap in a correctness argument (Section
\ref{sec:grid}). The new package \code{fastunknot} improves the exact part
substantially in practice (Section \ref{sec:fast}); its worst case is still
$2^{O(n)}$.

\section{The six archives}

\subsection{Extraction}
The six ZIP files in \code{reports/} were extracted into \code{reports/01}
through \code{reports/06}, each with the single wrapper directory removed, in
the order of the archive modification times:

\begin{center}\footnotesize
\begin{tabular}{@{}lllrr@{}}
\toprule
Dir & archive & package & Python lines & tests \\
\midrule
01 & \code{unknot\_recognition\_implementation.zip} & \code{unknot} (grid search) & 1753 & 40 \\
02 & \code{unknot\_recognition\_partial\_implementation.zip} & \code{unknotlab} & 1939 & 64 \\
03 & \code{unknot\_workbench.zip} & \code{unknot\_lab} & 1431 & 49 \\
04 & \code{unknot\_recognition\_partial.zip} & \code{unknot} & 1422 & 62 \\
05 & \code{unknot\_recognition\_partial\_implementation (1).zip} & \code{unknot\_recognition} & 1882 & 58 \\
06 & \code{unknot-recognition-implementation.zip} & \code{unknot} & 1278 & 45 \\
\bottomrule
\end{tabular}
\end{center}

Every archive contains a README, a LaTeX report with PDF, examples, recorded
results, and a standard-library-only Python package. All six test suites were
re-run here on CPython 3.14.4 under Windows (the archives recorded CPython
3.13.5 on Linux) and pass with exactly the claimed number of tests.

\subsection{What each archive delivers}

\paragraph{Archive 01: rectangular diagrams and Dynnikov moves.}
The only archive that does not use Khovanov homology. It represents knots as
grid diagrams, searches the finite space of grids reachable by
non-increasing moves (exchanges and destabilizations, modulo cyclic shifts)
and relies on Dynnikov's monotonic simplification theorem
\cite{dynnikov}: an unknot reaches the $2\times2$ grid, a knot is certified by
a state set closed under all moves. Certificates (positive path, non-unit
determinant, closed set) are replayed by a verifier that enumerates moves
independently. Bound $2^{O(g\log g)}$ in the grid size $g$; the report is
careful that $g$ is not the crossing number. It validated all 44\,719
connected grids of size $\le6$ (1285 shift orbits). It also contains the
pattern test, exact rational $H^1$ of a simplicial complex, cocycle-dual
normal coordinates, and the potential arithmetic.

\paragraph{Archive 02: Khovanov cube with fast certificates.}
Descending-diagram test (a diagram traversed meeting every crossing first on
its over-strand is trivial), Fox determinant, then the full reduced Khovanov
complex over $\F$. Its normal-surface module accepts genuine face-pairing
triangulations, checks orientability and sphere/disc vertex links, and
computes cocycle-dual normal coordinates without expanding discs. Six
braid/determinant pairs were checked against Wolfram \code{KnotData}.

\paragraph{Archive 03: workbench with an explicit exponential lower bound.}
Khovanov cube with default resource ceilings (65\,536 states), a
bond-enumeration pattern test with checkable witnesses, and a
\code{bound} command for the conditional $L(g+1)^L$ arithmetic. Its report
proves that its own implementation is not quasi-polynomial: on the unknot
$\widehat{\sigma_1\sigma_2\cdots\sigma_n}$ the cube has $2^n$ states and
$3^n$ generators, and the determinant filter cannot help since the
determinant is 1.

\paragraph{Archive 04: Khovanov cube with Knot Atlas fixtures.}
Khovanov cube, dual-4-connectivity pattern test, potential arithmetic. It
checks $10_{124}=T(3,5)$ from an Atlas PD code and contains a clean
chain-level proof that the unreduced $\F$ complex is two copies of the
reduced one, which justifies reading reduced ranks off published unreduced
tables.

\paragraph{Archive 05: Khovanov cube with Reidemeister preprocessing.}
Crossing-decreasing Reidemeister I/II moves with replayable traces,
determinant filter, Khovanov cube; bond-enumeration pattern test;
simplicial normal coordinates with manifold validation; the
pattern-complexity $-4\chi+|S\cap P|$ of the preprint. It exhausts all 682
two-strand words of odd length $\le9$ and gives an independent braid-theoretic
proof that its eight-crossing ``hard unknot'' is trivial.

\paragraph{Archive 06: Khovanov cube with 11-crossing tests.}
Khovanov cube, bond-enumeration pattern test, a \code{verify} command that
recomputes a report with $d^2=0$ checks. Its fixtures include the
Kinoshita--Terasaka and Conway knots (Alexander polynomial 1, reduced rank
33). Its source audit notes that a 34-page Oxford version of the slides uses
the factor $1/10$ instead of $1/3$ in the Cheeger condition and displays
$O((\log n)^2)$ for the hierarchy length; substituting that depth into the
digit count would give $n^{O((\log n)^2)}$.

\subsection{Common structure of the six reports}
Read side by side, the reports agree on every substantive point, which is
itself evidence since they were produced independently:
\begin{enumerate}
\item The requested bound is not established; the delivered recognizer is a
  different, exponential algorithm; resource limits return \code{UNKNOWN}
  and never a verdict.
\item The Khovanov route is justified by Kronheimer--Mrowka \cite{km} plus a
  universal-coefficient argument: $\dim_{\F}\Kh(K;\F)\ge\operatorname{rank}\Kh(K;\mathbb Z)$,
  and the rational reduced homology is nonzero because its Euler
  characteristic is the reduced Jones polynomial, so $\F$-rank one forces
  rational rank one.
\item Determinant one is never a positive certificate ($T(3,5)$ and the
  Kinoshita--Terasaka knot are the standard witnesses).
\item The talk's ``base $g$'' digit argument needs base $g+1$ when the digit
  $g$ is allowed, since $(1,0)$ and $(0,g)$ have the same base-$g$ value; all
  six use base $g+1$, as does Proposition 9.1 of the preprint.
\item The terminal pattern test is implemented only under an explicit
  ``already a 3-ball'' hypothesis, and $b_1=0$ is not ball recognition
  ($S^2\times I$ is the standard counterexample).
\item The missing operations are listed by name: initial layered handle
  structure of the exterior, compressed cutting with provenance,
  hierarchical multi-surfaces with the pattern-sensitive quadratic bound,
  lifting a violating dual cycle to a disc in an earlier stage, Cheeger-region
  detection and Heegaard modification, Haken's lemma and weak reduction, and
  an end-to-end bit-cost proof including restarts.
\end{enumerate}

\section{Code review}

\subsection{The Khovanov recognizers (02--06)}
All five have the same architecture, evidently converged upon rather than
copied (variable names, data layouts and gradings differ): enumerate the
$2^n$ resolutions with a disjoint-set structure, encode a reduced generator as
a bit mask over the unmarked circles with the marked circle fixed to $x$,
build each differential column as a Python integer (XOR arithmetic over
$\F$), and compute ranks by pivot elimination. Points checked in the review:
\begin{itemize}
\item Input validation is uniformly careful: one component, spherical rotation
  system via $V-E+F=2$ on the ribbon graph (so virtual diagrams are rejected,
  not misclassified), labels normalized, \code{\{"pd": []\}} defined as one circle.
\item The Frobenius algebra tables ($m(x\otimes x)=0$, $\Delta(1)=1\otimes x+x\otimes1$,
  $\Delta(x)=x\otimes x$) are identical everywhere, and every implementation
  checks that the marked circle keeps its $x$ label.
\item $d^2=0$ is available as an exhaustive optional check in all five and
  is exercised by their tests.
\item Gradings are deliberately unnormalized; only the total rank is used, so
  writhe shifts are irrelevant. This is correct for recognition.
\item Resource semantics are consistent (\code{UNKNOWN}, exit code 2 or 3;
  the codes differ between archives, which is a minor interoperability wart).
\item Archive 03 imposes default ceilings, so its default invocation is
  incomplete above 16 crossings; the README says so.
\end{itemize}
We found no correctness problem. The determinant filters (Fox colouring
matrix, Bareiss elimination with exact-division checks) are also equivalent.

\subsection{The pattern testers (all six)}
Three different criteria are implemented for ``is this cubic spherical
pattern essential on a ball'': dual simple and 4-connected with the $K_3$,
$K_4$ exceptions (01, 04), no dual loop, parallel pair or non-facial dual
triangle (02), and no bond of size 1 or 2 and no three-edge bond with more
than one vertex on each side (03, 05, 06). Each report proves its own
criterion from the definition. The proofs are correct; the bond argument in
archive 06 (the cycle rank of a side with $k$ vertices and $r$ cut edges is
$(k-r+2)/2$) is the cleanest. Input conventions differ (darts $2e,2e+1$;
signed edge ends; positional darts; edge identifiers occurring twice), which
is why a converter was needed for cross-validation.

\subsection{Normal coordinates and potentials}
Archives 01, 02 and 05 compute integral cocycles spanning $H^1(\,\cdot\,;\mathbb Q)$
in a spanning-tree gauge and convert a cocycle to normal coordinates by
half-integral level sets of local vertex heights. The constructions agree
(seven coordinates per tetrahedron, matching equations checked, Euler
characteristic without disc expansion). All are honest that this is not a
saturated $\mathbb Z$-basis, not a connected surface, and not a cut. The four
potential modules implement the same base-$(g+1)$ lemma and differ only in
interface.

\subsection{The grid search of archive 01}
\label{sec:grid}
This is the one place where the review found a gap in an argument. The
recognizer's negative certificate is a set of grids closed under all
implemented moves; soundness needs the implemented move set to be exactly a
set for which monotonic simplification holds. The report states that
``rows sharing an endpoint are not accepted by this implementation's
exchange rule'' and cites Dynnikov's Proposition 6, but does not argue that
Dynnikov's theorem still holds when exchanges of two adjacent rows (or
columns) with a common corner column are excluded. Depending on the reading
of Dynnikov's definition of interleaving, that exclusion may remove legal
moves, in which case a closed set could exist for an unknot. We could not
settle this from the sources at hand, so we tested it: Section
\ref{sec:xval-grid} compares the search verdict with Khovanov homology on
every grid of size $\le6$ and on random larger grids, and found no
disagreement. The gap is in the proof, not (as far as tested) in the answers.

Everything else in archive 01 checked out: the move implementation, the
Booth canonicalization, the independent verifier, and the exhaustive counts
($g!(g-1)!/2$ connected grids per size).

\section{Cross-validation}
None of the archives compared itself with another; each ran its own tests.
We wrote three scripts (in \code{synthesis/data/}) that load the packages
side by side.

\subsection{Five Khovanov implementations on shared inputs}
The same braid words and Knot Atlas PD codes were passed to all five packages
(after converting label conventions). Reduced $\F$ ranks:

\begin{center}\small
\begin{tabular}{@{}lrrrrrl@{}}
\toprule
Input & 02 & 03 & 04 & 05 & 06 & agree \\
\midrule
unknot circle & 1 & 1 & 1 & 1 & 1 & yes \\
trefoil & 3 & 3 & 3 & 3 & 3 & yes \\
figure-eight & 5 & 5 & 5 & 5 & 5 & yes \\
T(2,5) & 5 & 5 & 5 & 5 & 5 & yes \\
T(2,7) & 7 & 7 & 7 & 7 & 7 & yes \\
T(3,4) & 5 & 5 & 5 & 5 & 5 & yes \\
T(3,5) det 1 & 7 & 7 & 7 & 7 & 7 & yes \\
unknot s1 s2 & 1 & 1 & 1 & 1 & 1 & yes \\
unknot s1s2s3s4 & 1 & 1 & 1 & 1 & 1 & yes \\
conjugated unknot & 1 & 1 & 1 & 1 & 1 & yes \\
hard unknot 05 & 1 & 1 & 1 & 1 & 1 & yes \\
KT 11n42 (Atlas PD) & 33 & 33 & 33 & 33 & 33 & yes \\
Conway 11n34 (Atlas PD) & 33 & 33 & 33 & 33 & 33 & yes \\
10\_124 (Atlas PD) & 7 & 7 & 7 & 7 & 7 & yes \\
\midrule
30 random braid closures (3--4 strands, 6--10 letters) & \multicolumn{5}{c}{ranks 1--13} & 30/30 \\
\bottomrule
\end{tabular}

\end{center}

All 44 inputs agree, including the two Alexander-polynomial-one 11-crossing
knots and 30 random closures whose ranks range from 1 to 13.

\subsection{Six pattern testers on random spherical patterns}
Random perfect matchings of the darts of $V$ trivalent vertices were kept
when every component's rotation system is spherical, converted to the six
input formats, and classified by all six testers (400 distinct systems per
vertex count) together with six named patterns:

\begin{center}\small
\begin{tabular}{@{}rrrrr@{}}
\toprule
Vertices & spherical systems tested & random matchings drawn & essential & disagreements \\
\midrule
2 & 400 & 883 & 102 & 0 \\
4 & 400 & 737 & 12 & 0 \\
6 & 400 & 1179 & 0 & 0 \\
8 & 400 & 2603 & 0 & 0 \\
10 & 400 & 6144 & 0 & 0 \\
\bottomrule
\end{tabular}
\\[6pt]
\begin{tabular}{@{}lll@{}}
\toprule
Pattern & verdict & all six agree \\
\midrule
theta & essential & yes \\
tetrahedron & essential & yes \\
cube & essential & yes \\
prism & inessential & yes \\
dodecahedron & essential & yes \\
dumbbell & inessential & yes \\
\bottomrule
\end{tabular}

\end{center}

No disagreement among 2006 patterns and three different criteria. (Random
matchings on many vertices are almost always inessential because short bonds
are common; the named cube and dodecahedron are the essential large examples.)

\subsection{Grid search versus Khovanov homology}
\label{sec:xval-grid}
We wrote an independent grid-to-PD converter (traversal of the grid, crossing
signs from the directions of the vertical and horizontal strands) and compared
archive 01's search (determinant filter disabled) with archive 04's Khovanov
rank:

\begin{center}\small
\IfFileExists{tables/grid_xval.tex}{\begin{tabular}{@{}rlrrrr@{}}
\toprule
Grid size & sample & diagrams & unknots & knots & disagreements \\
\midrule
2 & exhaustive (all shift orbits) & 1 & 1 & 0 & 0 \\
3 & exhaustive (all shift orbits) & 2 & 2 & 0 & 0 \\
4 & exhaustive (all shift orbits) & 6 & 6 & 0 & 0 \\
5 & exhaustive (all shift orbits) & 64 & 62 & 2 & 0 \\
6 & exhaustive (all shift orbits) & 1212 & 1185 & 27 & 0 \\
7 & random, $\le 14$ crossings & 300 & 286 & 14 & 0 \\
8 & random, $\le 14$ crossings & 200 & 188 & 12 & 0 \\
9 & random, $\le 14$ crossings & 100 & 82 & 18 & 0 \\
\bottomrule
\end{tabular}
}{(table pending)}
\end{center}

The knotted orbits of size 5 are the two trefoils (grid number 5) and those of
size 6 are 26 trefoil presentations and one figure-eight, exactly as the
grid-number theory predicts. No disagreement was found, which supports but
does not prove the completeness of the restricted move set. The sample stops
at size 9 because the search itself becomes expensive there (one size-9 grid
needed 900\,863 states, and the 100 grids took 1.9 hours), while the Khovanov
side stayed below 14 crossings by construction.

\subsection{What the cross-validation does and does not show}
Agreement of five independently written cubes is strong evidence that all
five compute the same invariant, and agreement with the published ranks
(trefoil 3, figure-eight 5, $T(3,5)$ 7, $K11n34$ and $K11n42$ 33) ties that
invariant to reduced Khovanov homology. It says nothing about the
quasi-polynomial question, and all runs are on diagrams of at most 14
crossings because the cubes cannot go further.

\section{What is missing for \texorpdfstring{$\qp$}{n\^{}O(log n)}}
\label{sec:missing}
Consolidating the six audits with our own reading of the sources, an
implementation with a proved $\qp$ bound would require the following, none
of which exists in any archive or, to our knowledge, in print.

\begin{enumerate}
\item \textbf{A representation.} A layered handle structure of the knot
  exterior with a Morse function, boundary-pattern labels, the ambient
  embedding in $S^3$ (used by the generalised Seifert construction of
  pp.~75--78), compressed normal-surface data, and regluing maps, with a
  proved bound on the bit size of all of this after arbitrarily many
  operations. The archives' normal-coordinate modules are local ingredients
  only.
\item \textbf{Bounded multi-surfaces.} An algorithm producing independent
  relative homology classes with embedded representatives whose
  \emph{pattern-sensitive} complexity is $O(n^2)$, and a proof that
  compression and pattern-compression preserve the invariants that allow
  later surfaces to be kept. A rational $H^1$ basis is not this.
\item \textbf{Compressed cutting.} Cutting, component extraction, Euler
  characteristic, orientation and pattern transport on compressed data
  without expanding parallel discs, with provenance sufficient to lift a
  violating dual cycle in a terminal ball to a disc in the correct earlier
  stage. Agol--Hass--Thurston orbit counting is a tool, not this procedure.
\item \textbf{Cheeger regions and weak reduction.} Recognition of the
  Heegaard--Morse hypotheses, construction of the modified generalised
  splitting, Haken's lemma and weak reduction, and a global progress measure
  bounding restarts. The inequality $3g(\partial M_i)\le\min\{g(H_i),g(H)-g(H_i)\}$
  is a test on three integers, not the operation.
\item \textbf{The depth bound.} A proof that the executable process, after
  every rebuild and restart, has hierarchy length $O(\log n)$ (the talk) or
  $O((\log n)^2)$ (the Oxford version cited by archive 06), together with
  polynomial digit bounds. The preprint's Section 10 gives a linear bound in
  the initial $q$-complexity, and its Section 9 explicitly declines to bound
  $g$ and $L$.
\item \textbf{The cost accounting.} Only after all of the above does the
  conditional lemma proved in every archive, $L(g+1)^L$ iterations times a
  polynomial cost per iteration times polynomially many restarts, become a
  running-time theorem: $\log\bigl((g+1)^L\bigr)=O((\log n)^2)$.
\end{enumerate}

The archives also list, and we agree with, the substitutions that do not
work: an exponential search with a time cap is not a total algorithm; a
polynomial iteration count does not bound work per iteration; expanding a
compressed normal surface loses the compression; $b_1=0$ is not ball
recognition; a Jones or Alexander test does not detect the unknot.

\section{A faster exact recognizer: \code{fastunknot}}
\label{sec:fast}

\subsection{Design}
Since the geometric route cannot be completed from the available sources,
the second task was addressed by making the \emph{exact} route as fast as we
could while keeping every verdict provable. The package \code{fast/fastunknot}
(standard library only, Python $\ge3.10$) runs:
\begin{enumerate}
\item validation (PD codes, braid closures, or grid diagrams in the archive-01
  convention, converted by the traversal of Section \ref{sec:xval-grid});
\item crossing-decreasing Reidemeister I/II moves, as in archive 05;
\item the descending-diagram test of archive 02;
\item the \emph{Alexander polynomial}, computed exactly by fraction-free
  Bareiss elimination over $\mathbb Z[t]$ on a first minor of the Fox matrix
  of the Wirtinger presentation (rows $(1-t)x_o+tx_u-x_v$ at positive and
  $(1-t)x_o-x_u+tx_v$ at negative crossings, with signs from the traversal);
  $\Delta\ne1$ certifies knottedness and strictly dominates the determinant
  filter of the archives ($T(3,5)$ and all torus knots are now rejected in
  polynomial time; the Kinoshita--Terasaka and Conway knots are not);
\item reduced Khovanov rank over $\F$ by \emph{scanning}.
\end{enumerate}

\subsection{The scanning computation}
This is Bar-Natan's local algorithm \cite{barnatan} specialised to $\F$ and
implemented directly in terms of the arc-algebra description of morphisms.
The crossings are processed one at a time in a greedy order chosen (from up to
twelve starting crossings) to keep the boundary of the processed region small.
After $k$ crossings the partial complex has:
\begin{itemize}
\item objects: perfect matchings of the boundary edges, each with a
  homological degree; closed circles are delooped, one object with a circle
  becoming two objects (labels $1$ and $x$);
\item morphisms $a\to b$: $\F$-combinations of dot-subsets of the circles of
  $a\cup\bar b$, i.e.\ elements of $A^{\otimes\text{circles}}$ with
  $A=\F[x]/(x^2)$.
\end{itemize}
Composition is the two-dimensional TQFT of the glued surface: a component of
genus $\ge1$ or with two dots is zero (because the handle operator
$m\circ\Delta=2x$ and $x^2$ vanish), a closed component is $1$ exactly when it
carries one dot, and a planar component with $\beta$ boundary circles and
$d\in\{0,1\}$ dots is $\Delta^{(\beta-1)}(x^d)$, which over $\F$ is
$x\otimes\cdots\otimes x$ for $d=1$ and the sum over one undotted position for
$d=0$. Delooping is the same computation with a cup on the source circle
(dotted for label $x$) and a cap on the target circle (dotted to extract the
coefficient of $1$, since the dual basis of $\{1,x\}$ under the counit is
$\{\varepsilon(x\,\cdot\,),\varepsilon\}$). After every crossing all invertible
entries are cancelled: an entry $b\to c$ between equal matchings is a unit of
the local ring $\operatorname{End}(a)\cong A^{\otimes k}$ iff its expansion
contains the undotted term, its inverse is the finite geometric series of the
nilpotent part, and Bar-Natan's cancellation lemma replaces
$\varepsilon$ by $\varepsilon-\gamma\varphi^{-1}\delta$. When the last crossing is
added the boundary is empty, every morphism is a scalar, so a minimal complex
has zero differential and its number of objects is the unreduced total rank;
the reduced rank is half of it (archive 04's two-copies lemma) and equals 1
exactly for the unknot.

\begin{proposition}[Correctness]
On a validated one-component diagram, \code{fastunknot} returns \code{UNKNOT}
iff the knot is trivial, unless a user-selected ceiling stops it, in which
case it returns \code{UNKNOWN}.
\end{proposition}
\begin{proof}
Steps 2--3 are isotopies and a sufficient triviality criterion; step 4 uses
$\Delta_{\text{unknot}}=1$. Step 5 computes the homology of the Khovanov
complex: the tensor product of crossing complexes in Bar-Natan's category,
delooping, and Gaussian elimination are chain homotopy equivalences, and the
closed complex is the standard one. The decision rule is the detection theorem
with the universal-coefficient argument of the archives.
\end{proof}

\begin{proposition}[Cost]
Steps 1--4 are polynomial in $n$. Step 5 is $2^{O(n)}$ in the worst case.
\end{proposition}
\begin{proof}
The Alexander determinant is $O(n^3)$ polynomial operations of degree
$\le n$ with polynomially bounded coefficients. In step 5 the objects before
elimination after crossing $k$ number at most $2\cdot2^{c}$ times the objects
after crossing $k-1$, where $c\le2$ closed circles can appear, so the whole
computation is dominated by the cube; compositions cost polynomial time in the
boundary size.
\end{proof}

The interesting quantity is the size after elimination. Empirically it stays
close to (number of crossingless matchings of the boundary) $\times$ (a
multiplicity that grows with the tangle's homology), i.e.\ roughly
$2^{O(w)}\cdot\mathrm{poly}$ for scan width $w$. \textbf{We have no proof of
this}, and there cannot be one in general for the multiplicity: the reduced rank
of $\#_k\,K11n34$ is $33^k$ with Alexander polynomial 1, so any method that
outputs the full homology is exponential on such inputs, whatever the scan
width. A subexponential unknot recognizer along these lines would have to
detect ``rank $>1$'' without finishing the computation; we do not know how.

\subsection{Validation}
The package has 19 unit tests (polynomial arithmetic; Alexander polynomials
of the trefoil, figure-eight, $5_2$, $T(3,5)$, $\Delta=1$ for the two
11-crossing knots and two hard unknots; symmetry $\Delta(t)=\Delta(t^{-1})$,
$|\Delta(1)|=1$, mirror invariance and odd determinant on 40 random braids;
Reidemeister moves; the local endomorphism ring; $d^2=0$ after every
crossing; independence of the processing order; resource-limit semantics;
known ranks; an independent dense reduced-cube computation on 25 random
braids; the command line). Beyond the tests, the scanner agrees with archive
04's cube on 120 random braid closures (3--5 strands, 6--13 letters, with
random processing orders), and reproduces the Atlas ranks 33, 33 and 7.

\subsection{Measurements}
Timings are on the same Windows machine as everything else in this report;
they illustrate the dependence on scan width and are not complexity
evidence.

\begin{center}\small
\IfFileExists{tables/benchmark.tex}{\begin{tabular}{@{}lrrrrrr@{}}
\toprule
Family & $n$ & reduced rank & max.\ boundary & objects before & objects after & seconds \\
\midrule
unknot sigma\_1..sigma\_n & 8 & 1 & 2 & 6 & 2 & 0.001 \\
unknot sigma\_1..sigma\_n & 16 & 1 & 2 & 6 & 2 & 0.004 \\
unknot sigma\_1..sigma\_n & 32 & 1 & 2 & 6 & 2 & 0.013 \\
unknot sigma\_1..sigma\_n & 64 & 1 & 2 & 6 & 2 & 0.036 \\
unknot sigma\_1..sigma\_n & 128 & 1 & 2 & 6 & 2 & 0.133 \\
unknot sigma\_1..sigma\_n & 256 & 1 & 2 & 6 & 2 & 0.475 \\
T(2,n) & 5 & 5 & 4 & 30 & 10 & 0.003 \\
T(2,n) & 11 & 11 & 4 & 66 & 22 & 0.011 \\
T(2,n) & 21 & 21 & 4 & 126 & 42 & 0.043 \\
T(2,n) & 41 & 41 & 4 & 246 & 82 & 0.153 \\
T(3,n) & 8 & 5 & 6 & 42 & 10 & 0.010 \\
T(3,n) & 10 & 7 & 6 & 54 & 14 & 0.022 \\
T(3,n) & 14 & 9 & 6 & 78 & 18 & 0.032 \\
T(3,n) & 16 & 11 & 6 & 90 & 22 & 0.045 \\
T(3,n) & 20 & 13 & 6 & 114 & 26 & 0.092 \\
T(3,n) & 22 & 15 & 6 & 126 & 30 & 0.089 \\
random 3-braid & 20 & 75 & 6 & 762 & 150 & 0.217 \\
random 3-braid & 40 & 321 & 6 & 3354 & 642 & 1.035 \\
random 4-braid & 15 & 5 & 6 & 30 & 10 & 0.015 \\
random 4-braid & 21 & 25 & 8 & 438 & 146 & 0.576 \\
random 4-braid & 31 & 45 & 8 & 604 & 216 & 0.947 \\
random 4-braid & 41 & 109 & 8 & 654 & 218 & 2.683 \\
random 5-braid & 24 & 13 & 6 & 109 & 42 & 0.099 \\
random 5-braid & 36 & timeout & timeout & timeout & timeout & 600.147 \\
random 6-braid & 31 & 85 & 8 & 548 & 216 & 1.194 \\
\bottomrule
\end{tabular}
}{(table pending)}
\end{center}

\begin{center}\small
\IfFileExists{tables/pipeline.tex}{\begin{tabular}{@{}lrllr@{}}
\toprule
Example & $n$ & verdict & deciding step & seconds \\
\midrule
conway.json & 11 & KNOTTED & reduced-khovanov-F2-scan & 0.0528 \\
figure\_eight.json & 4 & KNOTTED & alexander-polynomial & 0.0001 \\
grid\_determinant\_one\_knot.json & 10 & KNOTTED & alexander-polynomial & 0.0010 \\
grid\_scrambled\_unknot.json & 6 & UNKNOT & reidemeister-reduction & 0.0004 \\
hard\_unknot\_8.json & 8 & UNKNOT & reduced-khovanov-F2-scan & 0.0022 \\
kinoshita\_terasaka.json & 11 & KNOTTED & reduced-khovanov-F2-scan & 0.0593 \\
torus\_3\_5.json & 10 & KNOTTED & alexander-polynomial & 0.0009 \\
trefoil.json & 3 & KNOTTED & alexander-polynomial & 0.0001 \\
unknot.json & 0 & UNKNOT & reidemeister-reduction & 0.0000 \\
unknot\_braid40.json & 40 & UNKNOT & reidemeister-reduction & 0.0090 \\
\bottomrule
\end{tabular}
}{(table pending)}
\end{center}

One case timed out: a random 36-letter closure on five strands exceeded the
600-second budget, while a 31-letter closure on six strands took about a
second. Both scans have boundary at most ten edges; the difference is the
size of the minimal complexes of the intermediate tangles, the multiplicity
factor discussed above, which is exactly the unproved part of the cost.

Two comparisons with the archives' cubes on identical inputs: the
Kinoshita--Terasaka knot (11 crossings, 2048 states, 11\,559 generators) takes
0.24--0.31 s in the cubes and 0.07 s in the scan with a boundary of six edges
and at most 66 objects after elimination; the unknot
$\widehat{\sigma_1\cdots\sigma_n}$, archive 03's explicit exponential family,
needs $2^n$ states in every cube and is decided by the scan in 0.011 s for
$n=80$ with two objects after every elimination.

\subsection{Limits of the new package}
It is a reference implementation in pure Python; a compiled implementation of
the same algorithm (as in Bar-Natan's or the KnotJob programs) would be far
faster. Its scan-order heuristic is greedy; for diagrams whose Tait graph is
grid-like the width is $\Theta(\sqrt n)$ at best and the minimal complexes can
be large. It does not produce a spanning disc or a Reidemeister sequence. And,
to repeat, it is not quasi-polynomial.

\section{Assessment and recommendations}

\paragraph{What was achieved.}
Seven mutually consistent exact recognizers; a verified and cross-validated
body of code for the terminal ball-pattern test and for cocycle-dual normal
surfaces; a precise, sourced statement of the conditional
$L(g+1)^L$ accounting and of everything a $\qp$ implementation would still
need; and one recognizer whose practical reach is well beyond the
$2^n$-state cubes.

\paragraph{Obstacles that remain.}
The six items of Section \ref{sec:missing}. They are research problems in
3-manifold topology and algorithm design, not coding tasks: the specification
of the accelerated hierarchy exists only as a sketch on slides, and the one
detailed paper on the subject states that its steps are not specified
precisely enough for a running-time estimate. Until that changes, any
``$\qp$ implementation'' would be an exponential search with a mislabelled
bound, which is exactly what all six archives refused to deliver.

\paragraph{If work continues.}
(1) Track the published record: a full version of the 2021 announcement, or
the refined Section 10 algorithm with an explicit cost model, would change the
picture. (2) Settle the move-set question of Section \ref{sec:grid} from
Dynnikov's paper, or add shared-corner exchanges to archive 01 and re-run its
exhaustive validation. (3) Port the scanning recognizer to a compiled
language and replace the greedy order by a tree-decomposition-based one; the
planar separator theorem gives width $O(\sqrt n)$ for the projection graph,
which would make the empirical behaviour subexponential even if unproved.
(4) Use the archives' exact recognizers as oracles for any future geometric
implementation, as several of their reports suggest.

\begin{thebibliography}{9}\small
\bibitem{slides} Marc Lackenby, \emph{Unknot recognition in quasi-polynomial time}, slides, February 2021, 109 pages, supplied as \code{docs/quasipolynomial-talk.pdf}.
\bibitem{lackenby2026} Marc Lackenby, \emph{Incompressible surfaces, hierarchies and unknot recognition}, arXiv:2607.23350v1, 25 July 2026, 54 pages. \url{https://arxiv.org/abs/2607.23350}.
\bibitem{km} P.~B. Kronheimer and T.~S. Mrowka, \emph{Khovanov homology is an unknot-detector}, Publ.\ Math.\ IH\'ES 113 (2011), 97--208.
\bibitem{barnatan} Dror Bar-Natan, \emph{Khovanov's homology for tangles and cobordisms}, Geom.\ Topol.\ 9 (2005), 1443--1499; and \emph{Fast Khovanov homology computations}, J.\ Knot Theory Ramifications 16 (2007), 243--255.
\bibitem{dynnikov} Ivan Dynnikov, \emph{Arc-presentations of links: monotonic simplification}, Fund.\ Math.\ 190 (2006), 29--76, arXiv:math/0208153.
\bibitem{khovanov} Mikhail Khovanov, \emph{A categorification of the Jones polynomial}, Duke Math.\ J.\ 101 (2000), 359--426.
\end{thebibliography}

\appendix
\section{Files produced}
\begin{itemize}
\item \code{reports/01} -- \code{reports/06}: the extracted archives (flat, one level).
\item \code{synthesis/report.tex}, \code{synthesis/report.pdf}: this report;
  \code{synthesis/make\_tables.py} regenerates the tables from
  \code{synthesis/data/*.json}; the cross-validation scripts are in
  \code{synthesis/data/}.
\item \code{fast/}: the \code{fastunknot} package, tests, examples,
  \code{benchmark.py} and \code{results/benchmark.json}.
\end{itemize}
\end{document}
